% Generated draft: requires mathematical and source review.
\documentclass[11pt]{article}
\usepackage[margin=25mm]{geometry}
\usepackage{fontspec}
\setmainfont{Latin Modern Roman}
\usepackage{amsmath,amssymb,mathtools,unicode-math}
\setmathfont{Latin Modern Math}
\usepackage{enumitem,needspace,xurl,hyperref}
\setlength{\parindent}{0pt}
\setlength{\parskip}{5pt}
\setlength{\emergencystretch}{3em}
\newcommand{\N}{\mathbb{N}}
\newcommand{\Z}{\mathbb{Z}}
\newcommand{\Q}{\mathbb{Q}}
\newcommand{\R}{\mathbb{R}}
\newcommand{\C}{\mathbb{C}}
\newcommand{\F}{\mathbb{F}}
\newcommand{\A}{\mathbb{A}}
\newcommand{\PP}{\mathbb{P}}
\newcommand{\E}{\mathbb{E}}
\newcommand{\eps}{\varepsilon}
\newcommand{\dd}{\,\mathrm{d}}
\begin{document}
\section*{UnsolvedMath ID: 800273}
\textbf{Source problem code:} catalogue entry
\subsection{Definitions and mathematical statement}
Let \(\PP^2(\C)\) be the complex projective plane, and let \(f:\PP^2(\C)\to\PP^2(\C)\) be a holomorphic map of degree \(d\geq 2\), as in the source. Write \(f^n\) for its \(n\)-fold iterate. A subset \(K\subseteq\PP^2(\C)\) is totally invariant if
\[
f^{-1}(K)=K=f(K).
\]
Define the critical set and the closed postcritical set by
\[
C=\{x\in\PP^2(\C):\operatorname{rank}_{\C}D f_x<2\},
\qquad
P_f=\overline{\bigcup_{n=1}^{\infty}f^n(C)},
\]
where closure is taken in \(\PP^2(\C)\).

The source asks: if \(K\) is totally invariant and
\[
P_f\cap K=\varnothing,
\]
must \(f\) be hyperbolic on \(K\)?

\textbf{Unresolved qualifications.} The supplied record does not define ``hyperbolic on \(K\)'', and does not require \(K\) to be closed, compact, or nonempty. These hypotheses must not be silently added. For example, if ``hyperbolic'' is intended to mean uniformly expanding, a precise candidate conclusion is that there exist constants \(a>0\) and \(\lambda>1\) such that
\[
\|D(f^n)_x v\|_{\mathrm{FS}}
\geq a\lambda^n\|v\|_{\mathrm{FS}}
\quad\text{for every }x\in K,\ n\geq1,
\ v\in T_x\PP^2(\C),
\]
where \(\|\cdot\|_{\mathrm{FS}}\) is the Fubini--Study norm. This is an explicitly conditional interpretation, not a definition established by the supplied source; the intended meaning of hyperbolicity remains to be resolved.
\subsection{Short English statement}
For a holomorphic self-map of the complex projective plane of degree at least two, suppose a totally invariant set avoids the closure of all forward images of the critical set. Must the map be hyperbolic on that set? The supplied formulation leaves the meaning of hyperbolicity and the regularity assumptions on the invariant set unspecified.
\subsection{Sources}
\begin{itemize}
\item Supplied local UnsolvedMath registry record: permanent ID \(800273\), source code \texttt{catalogue entry}, title ``Several Complex Variables Problem List --- Problem 72.'' The registry attributes the problem to Buzzard, Fornaess, and Krantz, with proposed year 1995.
\item Buzzard, Fornaess, and Krantz, \emph{The Several Complex Variables Problem List} (1995), Problem 72. Supplied source URL: \url{https://arxiv.org/abs/math/9509202}. The source record reports access on 2026-10-16 and extraction by \texttt{arxiv\_latex}; these are supplied provenance claims, not independently verified here.
\item The supplied registry labels the problem ``open'' and records an update dated 2026-10-16. This is supplied registry status, not an independently verified assessment of the literature.
\end{itemize}
\end{document}
